\documentclass{article}
\usepackage{../math_template}

\author{Jonathan Kasongo}
\title{Ersatz Math Olymipad 2024 --- P1/6}

\begin{document}
\maketitle

\begin{problem}{Ersatz Math Olymipad 2024 --- P1/6}
There are 1001 stacks of coins $S_1, S_2, \ldots, S_{1001}$.
Initially, stack $S_k$ has $k$ coins for each $k=1,2, \ldots, 1001$.
In an operation, one selects an ordered pair $(i, j)$ of indices
$i$ and $j$ satisfying $1 \leq i<j \leq 1001$ subject to two conditions:

\begin{itemize}
\item The stacks $S_i$ and $S_j$ must each have at least one coin.
\item The ordered pair $(i, j)$ must not have been selected in any
previous operation.
\end{itemize}

Then, if $S_i$ and $S_j$ have $a$ coins and $b$ coins, respectively,
one removes $\gcd(a, b)$ coins from each stack. \vspace{.2cm}

What is the maximum number of times this operation could be performed?
\end{problem}

\begin{solution}{Solution}
The answer is $250500$ times.
Let $P(i,j)$ denote
performing the given operation on the pair $(i,j)$.
Define the $n$-\textit{shrink} process to be as follows. Let $S_n$ be the
stack with the most amount of coins.
Perform the operations

$$
P(n-1,n) \rightarrow P(n-2,n) \rightarrow P(n-3,n) \rightarrow \cdots
$$

until we have performed $P(1,n)$ or the next operation would involve using
a stack with zero coins.\vspace{0.2cm}

Let $a<b$. A $b$-shrink and an $a$-shrink never repeat operations because
the ordered pairs $(\text{\_}, b)$ and $(\text{\_}, a)$ are never the same
since $a < b$.
Suppose that the $n$-shrink process terminates
after the operation $P(k,n)$, so the stack $S_k$ has zero coins after the
$n$-shrink.
Then all of the stacks $S_{n-1}, S_{n-2},
\ldots, S_k$ have had one coin removed from them, and the stack $S_n$ has
had $n-k$ coins removed from it. Observe the following property of the
$n$-shrink

$$
(1,2,3,\ldots,n-2, n-1, n) \overset{n-\text{shrink}}{\rightarrow}
(0,1,2,\ldots,n-3, n-2, 1)
$$

that is an $n$-shrink performed on $1,2,\ldots,n$ always changes the last
number to 1 and the first number to 0. Also we observe that if
an $n$-shrink uses $x$ operations, then the following $(n-1)$-shrink
would use $x-2$ operations. \vspace{0.2cm}

Thus if we perform a $1001$-shrink followed by a $1000$-shrink followed and
so on up to a $503$-shrink the sequence $(S_n)$ will look like so
$
0, 0,\ldots 0, 1, 2, 3, 1, 1, \ldots, 1
$
where the zeros repeat 499 times and the ones repeat 499 times. This is
because the after the 503-shrink we have performed $1001-503+1=499$ shrinks.
Now after the 502-shrink the subsequence $1,2,3$ gets turned into
$0,1,1$. And so now the overall sequence consists of 500 zeros, followed
by 501 ones.
\vspace{0.2cm}

So in total we have performed 500 shrinks which means we have used the
given operation
$2+ 4 + \cdots + 1000 = 250500$ times. Now we just need to show that
250500 is in fact maximal.\vspace{0.2cm}

If an operation $P(i,j)$ has $i\leq 500$, then there can be at most
$1+2+\cdots +500 = 500\cdot 501/2$ of these operations since each operation
will decrease the number of coins in stack $S_i$ by at least 1.
If an operation $P(i,j)$ has $i > 500$ then there can be at most
$\binom{501}{2} = 500\cdot 501/2$ such operations since each pair $(i,j)$
can be used once. Thus the upper bound for operations is $500\cdot 501=
250500$ as desired. $\blacksquare$
\end{solution}

\end{document}
